%& /Users/ziky/.cache/org-persist/17/f85e4f-97ce-446f-b8f4-86226120c229-f2d5840721536a4146334151303d1f5e
% Created 2026-02-02 Mon 13:51
% Intended LaTeX compiler: pdflatex
\documentclass[11pt]{article}
\usepackage[utf8]{inputenc}
\usepackage[T1]{fontenc}
\usepackage{amsmath}
\usepackage{amssymb}
\usepackage{capt-of}
\usepackage{hyperref}
\setlength{\abovedisplayskip}{0pt}
\setlength{\belowdisplayskip}{0pt}
\usepackage[a4paper, margin=1in]{geometry}
\usepackage{amsmath}

%% ox-latex features:
%   !announce-start, !guess-pollyglossia, !guess-babel, !guess-inputenc,
%   maths, !announce-end.

\usepackage{amsmath}
\usepackage{amssymb}

%% end ox-latex features


% end precompiled preamble
\ifcsname endofdump\endcsname\endofdump\fi

\author{Ziky Zhang}
\date{\today}
\title{Feb 2nd Class Notes}
\hypersetup{
 pdfauthor={Ziky Zhang},
 pdftitle={Feb 2nd Class Notes},
 pdfkeywords={},
 pdfsubject={},
 pdfcreator={},
 pdflang={English}}
\begin{document}

\maketitle
\section{Outline}
\label{sec:org8cc34ef}
Schrodinger's Equation
\begin{itemize}
\item Motivation
\item The Equation (Time independent)
\begin{itemize}
\item Exploration
\item Interpretation
\end{itemize}
\item Examples
\begin{itemize}
\item Infinite well
\item Finite well
\item Harmonic Oscillator
\end{itemize}
\item Time Dependence
\end{itemize}
\section{Motivation}
\label{sec:org091cbfe}
Framework for solving problems
-> First-Principles
   e.g. Newton's Laws of Motion
\section{Schrodinger's Equation (Time Independent)}
\label{sec:orga19b706}
\[ E \Psi (x) = - \frac{\hbar}{2m} \frac{\partial^2}{\partial x^2} \Psi (x)\ + V(x) \Psi (x) \]
\(\Psi (x)\) is a wave function
\begin{itemize}
\item Backed by Experiments
\item Keeps de Broglie's wavelength \(\lambda\)
\end{itemize}

\(\frac{\partial }{\partial x} \frac{\partial }{\partial x} \Psi (x)\) models a particle of mass subject to forces associated with \(V(x)\)
Interpretation of \(\Psi(x)\): Has something to do with liekly hood to find it as functinon of where you look
\begin{itemize}
\item The "Wave" of the function is a description for the graph of probability vs. distance.
\end{itemize}
\(\Psi(x)\) in general is complex \(\mathrm{d Prob}(x) = \Psi * \Psi \mathrm{d}x\) -> \(\Psi\) must be normalized -> \(\int^{\infty}_{-\infty} \mathrm{d Prob} = 1 = \int^{\infty}_{- \infty} \Psi * \Psi \mathrm{d}x\)
\subsection{Dimension analysis}
\label{sec:orgac603e5}
\begin{align*}
z^{\*} \text{only over } \mathbb{R}. \quad
z &= a + ib \\
z^{*} &= a - ib \\
z^{*} z &= (a - ib) ( a + ib) \\
  &= a^2 + b^2
\end{align*}
\begin{align*}
\mathrm{Dim}[\mathrm{Prob}] &= 1 \\
\mathrm{Dim}[\Psi^{*}\Psi \mathrm{d}x ] &= 1 \\
\mathrm{Dim}[\Psi] &= \frac{1}{\sqrt{\mathrm{length}}}
\end{align*}
\begin{align*}
\text{Weighted averge}: 
\left<  x \right> &= \int^{\infty}_{\infty} \Psi^{*} \cdot \Psi \cdot x \mathrm{d}x \\
\left< f(x) \right> &= \int^{\infty}_{\infty} \Psi^{*} \cdot \Psi \cdot f(x) \mathrm{d}x \\
f(x) &= x^2
\end{align*}
\subsection{Aside}
\label{sec:org1c9b492}
\begin{align*}
K(x) &= \frac{1}{2} mv^2 \\
  &= \frac{p^2}{2m} \\
  &= \frac{\hbar^2k^2}{2m} \\ \\
\Psi(x) &= e^{e(kx - \omega t)} \\
\frac{\partial}{\partial x}\Psi(x) &= ike^{e(kx - \omega t)} \\
\frac{\partial^2}{\partial x^2}\Psi(x) &= (ik)^2 e^{e(kx - \omega t)} \\
\frac{\partial^2}{\partial x^2}\Psi(x) &= - k^2 e^{e(kx - \omega t)} \\
\frac{\partial^2}{\partial x^2}\Psi(x) &= - k^2 \Psi(x) \\ \\
K(x) \Psi(x) &= \frac{\hbar^2k^2}{2m} \Psi(x) \\
K(x) \Psi(x) &= -\frac{\hbar^2}{2m} \frac{\partial^2 \Psi}{\partial x^2} \\
\end{align*}
\section{Infinite Well}
\label{sec:orgd4134e1}
In a potential well where \(x = 0\) and \(x = L\) has infinite potential,
\begin{align*}
v(x) = \left\{ 
  \begin{array}{ c l l }
    \infty & \quad \text{if } x < 0 & (I) \\
    \infty & \quad \text{if } x > L & (II) \\
    0      & \quad \text{if } 0 \leq x \leq L & (III)
  \end{array}
\right \}
\end{align*}
the force a particle will experience at \(x = L\) is \(F_x = - \frac{\mathrm{d}V}{\mathrm{d}x}|_{x = L} = - \infty\)

\begin{align*}
(III): -\frac{\hbar^2}{2m} \frac{\partial^2 \Psi}{\partial x^2} &= E \Psi \\
\frac{\partial^2 \Psi}{\partial x^2} &= -\frac{2mE}{\hbar^2} \Psi \\
\text{if: } \Psi(x) = A \sin(\kappa x) + B \cos(\kappa x) \\
\frac{\partial}{\partial x}\Psi(x) = A\kappa \cos(\kappa x) - B\kappa \sin(\kappa x) \\
\frac{\partial^2}{\partial x^2}\Psi(x) = - [A\kappa^2 \sin(\kappa x) + B\kappa^2 \cos(\kappa x)] \\
\frac{\partial^2}{\partial x^2}\Psi(x) = - \kappa^2 \Psi(x)
\end{align*}
\[
    \Psi(0) = 0 \\
    \Psi(L) = 0 
\]

\begin{align*}
\Psi(x) = A_n \sin(\kappa_n x) \\
\kappa_n = \frac{n \pi}{L}, \quad n=1, 2, 3, \dots \\
E_n = - \kappa_n^2 \Psi(x)_n \\
\end{align*}
\textbf{NOTES TO BE ADDED}
\section{Finite Well}
\label{sec:org93a7845}
\begin{align*}
v(x) = \left\{ 
  \begin{array}{ c l l }
    V_0 & \quad \text{if } x < 0 & (I) \\
    0   & \quad \text{if } 0 < x < L & (II) \\
    V_0 & \quad \text{if } L < x & (III) \\
  \end{array}
\right \}
\end{align*}
\begin{align*}
\mathrm{I:}  &E \Psi (x) = - \frac{\hbar}{2m} \frac{\partial^2}{\partial x^2} \Psi (x)\ + V(x) \Psi (x) \\
\mathrm{II:}  &E \Psi (x) = - \frac{\hbar}{2m} \frac{\partial^2}{\partial x^2} \Psi (x)\ \\
\mathrm{III:}  &E \Psi (x) = - \frac{\hbar}{2m} \frac{\partial^2}{\partial x^2} \Psi (x)\ + V(x) \Psi (x)
\end{align*}
Solution to this problem involes solving trancendental algebra problems. (e.g. \(2x = 3 \sin(5x)\))
Has strongg denpendence for number of trapped electron quantum states on \(L\text{, } m\text{, } V_0\) 
\section{Harmonic Oscillartor}
\label{sec:orga64811c}
\end{document}
